\section{Related work}\label{sec:related}

The construction sits between two literatures that rarely meet: the decomposition of
proper scores, which supplies the object being estimated, and the comparison of
forecasters, which supplies the inferential target. We take those first, then the applied
literature that motivates the problem.

\subsection{Relation to classical proper-score decompositions}\label{sec:related-decomposition}

Conditional on a chosen grouping variable, a proper score splits into a calibration
(reliability) term and a resolution term \citep{murphy1973vector,degroot1983comparison}, made
conditioning-set explicit by \citet{brocker2009reliability}, whose treatment already covers
general strictly proper scores and multicategory outcomes.

Our $\Delta R$ is related to that resolution term but is not equal to it, and the distinction
matters before anything is built on it. The representation of a resolution-like quantity as a
\emph{covariance} between forecast and outcome indicator is due to \citet{yates1982external},
whose covariance decomposition of the Brier score is the direct ancestor of
Eq.~\eqref{eq:cov-id}. The classical resolution term of \citet{degroot1983comparison} is a
different object: it depends on a forecast only through the partition generated by its level
sets, so it is invariant under any strictly monotone relabelling of forecast values. A
covariance between forecast values and outcomes has no such invariance. The two coincide ---
$\Delta R_c$ equals twice a difference of classical resolution terms --- exactly when both
models are calibrated within each cell of $\phi$, and otherwise differ by a
reliability--resolution cross term that a monotone recalibration moves.
Section~\ref{sec:simulation} measures the gap: a pure recalibration of one model, which
cannot change either model's classical resolution at all, moves $\widehat{\Delta R}$ by
$-0.488$ and reverses the sign verdict of the audit in Section~\ref{sec:sggaudit}. We
therefore name $\Delta R$ for the covariance it is, and adopt confidence matching as
procedure rather than caveat.

What the classical decomposition does not supply is inference for a pairwise contrast. The
point estimates do combine: both models share the $\phi$-partition and hence the same
per-cell $n_c$, so the same $(n_c-1)/n_c$ attenuation applies to each model's separately
estimated in-sample resolution, and subtracting two naive plug-ins recovers exactly the
biased estimate of the contrast. Inference does not transfer. Variance estimators for the
single-model decomposition \citep{siegert2014variance} target one resolution term in
isolation and supply no covariance between it and a second model's estimate from the same
data. Rebuilding that from the joint law is routine rather than open, and the field is not
silent about it \citep{ferro2007comparing,delsole2014comparing,delong1988comparing}. What is
new is not that a contrast can be given a variance, but that it admits an \emph{exact
finite-sample decomposition} whose remainder is an antisymmetric U-statistic with a one-pass
influence function, so the split and its uncertainty come from one construction.

\subsection{Comparing two forecasters}\label{sec:related-dm}

Inference on a paired score differential is the classical territory of comparative
forecast evaluation: the Diebold--Mariano test \citep{diebold1995comparing} and
conditional predictive ability tests \citep{giacomini2006tests} ask whether one forecaster
beats another, optionally conditional on covariates. That is the closest existing relative
of a gain evaluated conditional on $\phi$, and
Corollary~\ref{cor:dm} makes the relation exact: at the trivial one-cell partition our
total-gain statistic and its interval \emph{are} a clustered Diebold--Mariano test,
corrected for cluster rather than serial dependence. The decomposition is therefore best
read as everything those tests already deliver plus a split that neither provides. A
structurally close precedent for comparing two correlated models through a paired
U-statistic exists outside forecast evaluation as well: \citet{delong1988comparing}
compare two correlated ROC curves' areas via a generalized U-statistic covariance
estimator, the same construction applied to an unconditional ranking statistic rather than
a declared-cell gain.

Two further literatures locate the same object from other directions. The transported
component is precisely what label-shift procedures manipulate post hoc: re-weighting a
classifier's outputs towards a new label distribution
\citep{saerens2002adjusting,lipton2018detecting} can add or remove transported gain without
touching case-level evidence, which Appendix~\ref{app:vgvisual} exploits as a
falsification test. And evaluating calibration within cells of a declared grouping,
uniformly over many candidate groupings, is the subject of multicalibration
\citep{hebert2018multicalibration}, which is the formal treatment of the multiplicity
question raised by the choice of $\phi$ (Section~\ref{sec:discussion}).

\subsection{Transport, permutation, and U-statistics}\label{sec:related-perm}

Strictly proper scoring rules reward honest predictive distributions
\citep{gneiting2007proper}, and permutation methods have long been used to assess classifier
performance \citep{ojala2010permutation}. Conditional permutation tests generally estimate
or assume a conditional distribution when the conditioning variables are continuous or
high-dimensional. The setting here is different in a way that buys exactness: a benchmark
supplies a \emph{discrete} grouping variable, and the target is not generic conditional
independence but one specific model-to-model score gain, so within-cell transport needs no
nuisance law at all. The remainder is a U-statistic of order two in the sense of
\citet{hoeffding1948class}, and its H\'ajek projection is what supplies the influence
function of Section~\ref{sec:inference}. The new element is the gain-contrast transport
identity and the attribution it licenses, not permutation or U-statistic theory by itself.

Two further connections are worth naming. The transported component is what label-shift
procedures manipulate post hoc: re-weighting a classifier's outputs towards a new label
distribution \citep{saerens2002adjusting,lipton2018detecting} can add or remove transported
gain without touching case-level evidence, which Section~\ref{sec:vgvisual} exploits as a
falsification test, and which is why Section~\ref{sec:decision}'s model-selection reversal
is a label-shift argument. And what an unrecorded grouping variable can hide is the subject
of omitted-variable sensitivity analysis \citep{rosenbaum2002observational};
Proposition~\ref{prop:sensitivity} adapts that style of argument to bound how far an
unrecorded covariate could move $\Delta R$.

Usable-information analyses measure information accessible to a chosen predictor family
\citep{xu2020theory}. Their central objects are model- or family-level quantities; ours is
narrower and more operational, attributing the score difference between two concrete frozen
models. A positive $\Delta R$ establishes nothing about causal or compositional structure.

\subsection{Prior-aware evaluation in vision}\label{sec:related-applied}

Dataset bias has long complicated comparisons between recognition systems
\citep{torralba2011unbiased,geirhos2020shortcut}. In scene graph generation a frequency
predictor of $p(\text{predicate}\mid\text{subject},\text{object})$ is surprisingly
competitive \citep{zellers2018neural}; mean recall and causal interventions were proposed to
reduce that dependence \citep{chen2019knowledge,tang2019learning,tang2020unbiased}, though
metric pathologies remain \citep{li2022rethinking} and the same diagnosis recurs in
open-vocabulary and knowledge-enhanced work
\citep{li2024pixels,zhang2024hiker,sun2024unbiased}. In visual question answering, GQA-OOD
stratifies by answer rarity within question groups \citep{kervadec2021roses} and VQA-CP
changes language priors between train and test \citep{agrawal2018dont}, though
changing-prior splits can themselves be exploited \citep{teney2020value}; the shortcut is
still targeted directly \citep{xu2025qirl,kuang2024debiasing}, and benchmark-level audits of
exploitable non-visual shortcuts have widened to multimodal evaluation generally
\citep{brown2025benchmark,zhou2025shortcut}. In long-tailed recognition a classifier can
raise aggregate accuracy by fitting training class frequencies more precisely without
discriminating better within a class \citep{ma2024geometric,duan2024longtailed}; the
canonical remedies rebalance the objective or the classifier
\citep{cui2019classbalanced,cao2019learning,menon2021long,kang2020decoupling}, with a
continuing stream of variants \citep{li2026pih2t,luo2025rebalancing}.

Reporting standards have responded to the shortcut: mean recall alongside recall, often
broken down by predicate frequency \citep{tang2020unbiased,li2022rethinking}, and zero-shot
recall \citep{lu2016visual} checking whether a gain survives on triplets absent from
training. These are effective prior-aware slicings, splits and stress tests of a
\emph{single} model's performance. What none provides is what we target: an exact
decomposition, with uncertainty, of the score difference between two fixed systems. A
stratified or zero-shot metric can show \emph{that} a gain is uneven across the prior; it
does not attribute the pairwise gain itself, on the full evaluation set, with a
finite-sample identity and inference. The remedies, equally, intervene on how a model is
trained rather than on how a trained pair is compared.
